\begin{claim}{F.16}{Proved}{A full procyclic ambiguity invisible to every proper forgetting}
Use $P_8$, the actual involution $\mathsf r$, and the
source image
\[
       H=\widehat\phi\bigl(\widehat{\pi_1(Y,s_0)}\bigr)
                         \ \leq P_8
\]
of F.12--15.  There is an explicitly constructed discrete
element $w\in P\Gamma_{0,8}$ such that
\[
\begin{split}
 Z_w:=\overline{\langle w\rangle}&\cong\widehat\Z,\\
 Z_w&\subset C_B(\mathsf r),\\
 q_I(Z_w)&=1
       \quad\text{for every proper labelled forgetting }q_I,\\
 Z_w\cap H&=1 .
\end{split}                                                     \tag{F.16.1}
\]
Here $q_I$ retains any proper subset $I$ of the eight
labels, with the usual induced homomorphism on full
pure mapping class completions.  In particular all
eight individual one-point forgettings kill $Z_w$.

The exact image-membership test ruled out by this
construction is the following pullback subgroup:
\[
 \mathcal T=
       \bigcap_{I\ {\rm proper}}
                     q_I^{-1}\bigl(q_I(H)\bigr).
                                                               \tag{F.16.2}
\]
Both inclusions
\[
 H\ \subsetneqq\ \mathcal T,\qquad
 H\cap C_{P_8}(\mathsf r)\
       \subsetneqq\
          \mathcal T\cap C_{P_8}(\mathsf r)
                                                               \tag{F.16.3}
\]
are strict, witnessed by every nonidentity element of
$Z_w$.  The same strictness holds when the intersection
uses only the eight individual forgettings, or any
subfamily of proper forgettings.
For each $h\in H$, the elements $hz$ with $z\in Z_w$
have exactly the same images as $h$ under all these
maps, and $h$ is their only member in $H$.
Within each nonempty correction coset of F.15.9,
right multiplication by $Z_w$ likewise preserves
centralization and every proper forgetting image.
The original source-membership equation and the
full residual centralizer remain unresolved.
\end{claim}

\begin{proof}
\emph{The actual six-point quotient and its labelled point-pushing group.}
Keep the two fixed points $a_+,a_-$ from F.15,
with $a_-=-a_+$ and $a_+^2=c_0$.
The specified quotient of $r(z)=c_0/z$ is
\[
    q(z)=\left(\frac{z-a_+}{z+a_+}\right)^2.
                                                               \tag{F.16.4}
\]
Indeed the coordinate $(z-a_+)/(z+a_+)$ changes
sign under $r$.  Label the two quotient branch
points $A=q(a_+)=0$, $D=q(a_-)=\infty$.
The four ordinary quotient labels and their
literal fibres on the eight marked points are
\[
\begin{array}{c|c|c}
 i&o_i&\text{marked fibre}\\ \hline
 0&1&\{0,\infty\}\\
 1&q(1)&\{v_0,v_2\}\\
 2&q(v_1(s_0))&\{v_1,v_3\}\\
 3&q(v_4(s_0))&\{v_4,v_5\}.
\end{array}                                                     \tag{F.16.5}
\]
These six quotient points are distinct: the four
ordinary values are the four different free
$r$-orbits, and none is a branch orbit.
This is a quotient of the actual configuration,
with the same choice of $a_+$ as F.15.

In the pure mapping class group of
$(\mathbf P^1;A,D,o_0,o_1,o_2,o_3)$,
forget the branch label $A$.  The actual Birman
kernel and point-pushing map are
\[
 \begin{split}
 L&=\pi_1\bigl(\mathbf P^1\setminus
                    \{D,o_0,o_1,o_2,o_3\},A\bigr)\cong F_4,\\
 \operatorname{push}_A:L&\hookrightarrow P\Gamma_{0,6}.
 \end{split}                                                    \tag{F.16.6}
\]
Choose a fixed geometric system of four disjoint
whiskers from $A$ to small circles about the four
$o_i$, with their interiors avoiding all five
other labels.  Orient the circles positively.
Their based loops $m_0,m_1,m_2,m_3$ are the
specified free basis of $L$; the omitted
$D$-meridian is the inverse of their ordered
product.  The based point-pushing map is induced
by moving $A$ along these loops and fixing
the other five labels.  Its inverse on the
Birman kernel is the track of $A$ in an isotopy
after that label has been forgotten.  The
aspherical Birman fibration in F.14 makes this
inverse independent of the isotopy, and gives
both the displayed injection and the full
induced profinite topology.

Let $\gamma_i$ be the boundary of a small regular
neighbourhood of the $i$th whisker and its two
ends.  It encloses exactly $A,o_i$.
Write
$T_i=\operatorname{push}_A(m_i)$.
This is the one-turn Dehn twist about $\gamma_i$,
with its sign fixed by this point-pushing
convention.  To see the representative directly,
move $A$ once round the small circle at $o_i$
and back along its whisker; outside the
regular neighbourhood the endpoint isotopy is
the identity, and on its boundary annulus it
is precisely that one-turn twist.

\emph{The pure lift, including its annulus degree.}
The double-cover monodromy is nontrivial around
each of $A,D$, and trivial around every $o_i$.
Therefore the inverse image of $\gamma_i$ is one
connected annulus covering its downstairs
annulus with degree two.  The disk it bounds
upstairs is the double cover of the disk
enclosing $A,o_i$, branched at its one point
$A$.  Its Euler characteristic is $2-1=1$,
and its boundary is connected, so it is a disk.
It contains exactly the two marked points in
row $i$ of (F.16.5), together with the unmarked
fixed point $a_+$, and no other marked point.

Use angular coordinates in turns on the two
annuli, so the covering map is
$\theta=2\vartheta$ modulo one.
Choose the radial coordinate $t\in[0,1]$
according to the preceding twist convention.
The square of the downstairs twist and its
pure lift are exactly
\[
 \begin{split}
  T_i^2:(\theta,t)&\longmapsto(\theta+2t,t),\\
  b_i:(\vartheta,t)&\longmapsto(\vartheta+t,t).
 \end{split}                                                    \tag{F.16.7}
\]
Both boundary circles are fixed pointwise by
the second map.  It extends by the identity
off the annulus.  The deck action there is
$\vartheta\mapsto\vartheta+\tfrac12$, which
commutes with (F.16.7).
Thus $b_i$ is a pure mapping class that
centralizes the actual $r$ and descends
to $T_i^2$.
It is the unique pure lift: the only other
lift differs by $r$, whose permutation is the
four nontrivial pairs (F.16.5).

The class $b_i$ is the one-turn twist about
the boundary of the displayed two-puncture
disk.  Forgetting either of those two points
leaves a disk with one marked point.
Its twist is isotopic to the identity
relative to that point and all exterior
labels: extend the untwisting isotopy
across its disk by rotations about that
remaining point.  Consequently
\[
 b_i\in C_{P\Gamma_{0,8}}(\mathsf r)
        \ \cap\
        \ker q_{p_i}\ \cap\ \ker q_{p'_i},
 \qquad
       \{p_i,p'_i\}=q^{-1}(o_i)\cap
                     \{0,\infty,v_0,\ldots,v_5\}.
                                                               \tag{F.16.8}
\]
All kernels here refer to the actual named
point-forgetting maps.

\emph{The nontrivial word and all eight forgettings.}
Take the commutator convention
$[g,h]=ghg^{-1}h^{-1}$ and define
\[
 \begin{split}
 w&=[[[b_0,b_1],b_2],b_3],\\
 \omega&=[[[m_0^2,m_1^2],m_2^2],m_3^2]\in L .
 \end{split}                                                    \tag{F.16.9}
\]
Every $b_i$ centralizes $\mathsf r$, so $w$
does.  Forgetting any one of the eight
points kills the corresponding $b_i$
by (F.16.8); substituting the identity
for any input of the displayed nested
commutator gives the identity.
This proves all eight individual
vanishings.  Any proper forgetting factors
through one of them, which proves all
the asserted vanishings after completion
as well.
Descent through the actual quotient gives
\[
             \operatorname{desc}(w)
                       =\operatorname{push}_A(\omega).
                                                               \tag{F.16.10}
\]
In particular this is an identity in the
specified Birman kernel, with its based
inverse in (F.16.6).

Here is an integral detector for $\omega$.
Let $U_5(\widehat\Z)$ be the group of
upper unitriangular five-by-five matrices,
and let $E_{ij}$ denote its usual matrix
units.  The free basis map has a unique
continuous extension
\[
 \widehat L\longrightarrow U_5(\widehat\Z),
       \qquad m_i\longmapsto I+E_{i+1,i+2}
                              \quad(0\leq i\leq3).
                                                               \tag{F.16.11}
\]
The four squares go to $I+2E_{i+1,i+2}$.
Successive multiplication in the actual
unitriangular group gives
\[
 \begin{split}
 [I+2E_{12},I+2E_{23}]&=I+4E_{13},\\
 [I+4E_{13},I+2E_{34}]&=I+8E_{14},\\
 [I+8E_{14},I+2E_{45}]&=I+16E_{15}.
 \end{split}                                                    \tag{F.16.12}
\]
Thus $\omega$ is nontrivial already
discretely.  For every $z\in\widehat\Z$,
its full profinite power has image
\[
                         I+16zE_{15}.                          \tag{F.16.13}
\]
Multiplication by sixteen is injective
on $\widehat\Z$; this detects every
nonzero primary coordinate, including
the two-primary coordinate.

\emph{The precise completion maps for this detector.}
Let $N_r$ be the discrete normalizer of
the specified $\langle r\rangle$ in the
eight-puncture mapping class group.
Its decorated quotient sequence F.7.5
is
\[
       1\longrightarrow\langle r\rangle
         \longrightarrow N_r
         \xrightarrow{\ \operatorname{desc}\ }Q
         \longrightarrow1 ,
                                                               \tag{F.16.14}
\]
where the two branch labels and four
ordinary labels of (F.16.5) are retained.
Every pure six-point mapping class
preserves the double-cover monodromy:
it fixes the peripheral homology classes,
and that monodromy is the specified
$C_2$-valued homology character.
Thus $P\Gamma_{0,6}$ is a finite-index
subgroup of this $Q$.
The Birman kernel $L$ in (F.16.6) has
its full induced topology in
$P\Gamma_{0,6}$ by F.14, and hence in
$Q$.  In particular the actual completed
maps give embeddings
\[
           \widehat L\hookrightarrow
             \widehat{P\Gamma}_{0,6}
                    \hookrightarrow\widehat Q .
                                                               \tag{F.16.15}
\]

Set $C_r=C_{P\Gamma_{0,8}}(\mathsf r)$.
It is a finite-index subgroup of $N_r$.
Its intersection with the finite deck
kernel is trivial, since that nonidentity
element permutes the marked points.
Consequently descent identifies $C_r$
with its finite-index image $Q_0$ in $Q$.
The inverse on $Q_0$ is the unique
lift fixing all eight labels individually.
The prime-order normalizer and full-topology
comparison F.7.4--5 make the completion
of this isomorphism the actual continuous
descent on the closure of $C_r$ in $P_8$:
\[
       \overline{C_r}\cong\widehat{C_r}
             \xrightarrow[\ \operatorname{desc}\ ]{\ \sim\ }
                           \widehat{Q_0}
               \hookrightarrow\widehat Q .
                                                               \tag{F.16.16}
\]
In particular (F.16.10) persists under
these maps, and its right side lies in
the embedded $\widehat L$ of (F.16.15).
If $w^z=1$ in $P_8$, its descent is
$\omega^z=1$ in $\widehat Q$.
The embedding (F.16.15) then puts this
equality in $\widehat L$, and
(F.16.13) gives $16z=0$, hence $z=0$.
The usual continuous power map
$\widehat\Z\to\overline{\langle w\rangle}$
is therefore injective as well as onto.
This proves the full procyclic
identification in (F.16.1), with the
inverse-source and topology comparisons
needed to apply the matrix detector.

\emph{The exact obstruction to the proposed membership test.}
The eight vanishing maps include
$q_{v_0}$ and $q_{v_2}$, so
$Z_w\subset B$.  Its elements centralize
$\mathsf r$, proving $Z_w\subset C_B(\mathsf r)$.
F.14 states that $q_{v_0}|_H$ is
injective.  Since $q_{v_0}(Z_w)=1$,
we obtain $Z_w\cap H=1$.

For every proper $I$ the identity lies in
$q_I(H)$, and $q_I(Z_w)=1$.
Thus $Z_w\subset\mathcal T$, while
$H\subset\mathcal T$ by its definition.
Every nonidentity element of $Z_w$
witnesses the first strict inclusion.
It also centralizes $\mathsf r$ and
therefore witnesses the second.
The same witnesses work for every
subfamily of the displayed maps.
For $h\in H$ and $z\in Z_w$,
$q_I(hz)=q_I(h)$ for all $I$, whereas
$hz\in H$ forces $z\in H\cap Z_w=1$.
Finally $Z_w\subset C_B(\mathsf r)$
places right multiplication by $Z_w$
inside each actual centralizing
correction coset (F.15.9), while all
proper forgetting maps remain unchanged.
These statements prove precisely
(F.16.2)--(F.16.3); they supply no
nonidentity element of the original
source equalizer.

The executed certificate
\path{code/part_a/common_reflection_forgetting_ambiguity_f16.py}
records the four actual orbit pairs,
checks the degree-two annulus formulas,
all eight individual and all proper
word-forgetting consequences, the
integer matrix value (F.16.12), and
its power formulas in twenty-one
finite primary quotients.
The geometric lifting and full
completion maps are proved above.
\end{proof}
